\section{Nghiên cứu người dùng (User Research)}

\subsection{Target Users}
% Trình bày:
% - Người dùng mục tiêu (Chủ nuôi thú cưng bận rộn)
% - Đặc điểm nhân khẩu học và thói quen chăm sóc thú cưng
% - Context of use

\subsection{Research Methods}
% Trình bày phương pháp thực hiện (Interview, Survey, Observation...):
% - Mục đích
% - Số lượng và tiêu chí người tham gia
% - Quy trình thực hiện
% - Dữ liệu thu thập

\subsection{Research Findings}
% Trình bày các phát hiện và dữ liệu thực tế:
% - Người dùng thường gặp khó khăn khi...
% - Người dùng kỳ vọng...
% - Người dùng nhầm lẫn hoặc lo lắng về...

\subsection{User Needs / Pain Points}
% Tổng hợp nhu cầu và điểm đau chính của người dùng:
% - Pain Points: Chờ xác nhận lâu, thất lạc dặn dò đặc biệt, thiếu minh bạch tiến độ...
% - User Needs: Đặt lịch tức thì, tự động lưu hồ sơ dị ứng/thuốc, theo dõi tiến độ thời gian thực...

\subsection{Personas}
% Trình bày Persona dựa trên dữ liệu phỏng vấn thực tế:
% - Primary Persona: Nguyễn Mai Anh (28 tuổi, Nhân viên Marketing)
% - Goals, Behaviors, Pain points, Needs
